\justifying
\NSFCBodyText

（1）自治区高校基本科研业务费科研项目（培育类），项目编号 XJEDU2026P005，项目名称“融合多源数据大模型驱动的光伏发电功率预测方法研究”，经费 5 万元（直接经费 3.9 万元、间接经费 1.1 万元），执行期 2026 年 1 月至 2026 年 12 月，申请人主持，负责总体方案、模型研究和成果组织。该项目研究对象为光伏功率预测，与本申请的中文预训练语料质量机制在研究对象、数据、科学问题和成果上不同；其关联仅为深度模型研究与项目组织经验，不与本申请重复资助。

（2）新疆数字云创有限公司委托横向课题，合同编号 202604140013，项目名称“面向政企数字化应用的智能知识库与问答系统研发”，经费 12 万元，执行期 2026 年 6 月至 2026 年 10 月，申请人主持，合同任务为构建一个知识库，项目进行中，成果及知识产权由合同双方共同享有。该课题面向政企知识库工程交付，本申请面向中文预训练语料筛选效应的基础机制，两者不共用研究目标和考核成果；可复用的仅是数据治理和软件工程经验，任何合同数据与成果均不自动进入本申请。

上述项目的经费、数据和设备不计入本项目可支配资源。申请人将通过独立的数据版本、代码仓库、运行清单和成果标注管理项目边界，避免任务与成果重复归属。
